\section{Register Umum dan Register Segmen}
\label{sec:register-umum}

Seluruh pekerjaan prosesor dilakukan di atas sekumpulan tempat
penyimpanan kecil yang berada langsung di dalam cip prosesor. Tempat
penyimpanan itu disebut \term{register}. Register adalah memori
tercepat yang dimiliki komputer: satu register dapat dibaca atau
ditulis dalam satu daur mesin, jauh lebih cepat daripada mengakses
memori utama yang harus melewati bus.

Prosesor 8086 memiliki empat belas register yang dapat dipakai
programmer. Masing-masing berukuran 16 bit, sehingga dapat menampung
nilai \hex{0000} sampai \hex{FFFF}. Keempat belas register tersebut
dikelompokkan menjadi empat kelompok: register umum, register segmen,
register \eng{pointer} dan \eng{index}, serta register \eng{flag}.
Bab ini membahas keempat kelompok tersebut satu per satu.
Pembahasan mengenai organisasi internal prosesor secara menyeluruh
telah disinggung pada Bab~\ref{chap:arsitektur-8086}.

\subsection{Register Umum dan Akses Byte}

Empat register pertama, yaitu \reg{AX}, \reg{BX}, \reg{CX}, dan
\reg{DX}, disebut \term{register umum} (\eng{general purpose
register}). Disebut umum karena keempatnya dapat dipakai untuk
hampir semua keperluan: menampung hasil hitungan, menyimpan alamat,
menjadi penghitung, atau sekadar tempat sementara.

Keistimewaan keempat register ini terletak pada kemampuannya diakses
sebagai dua bagian 8 bit yang terpisah. Setiap register 16 bit
terbagi menjadi \term{byte tinggi} (\eng{high byte}) dan \term{byte
rendah} (\eng{low byte}), seperti tampak pada Tabel~\ref{tab:register-umum}.

\begin{table}[H]
    \centering
    \caption{Register umum 8086 dan akses byte rendah/tinggi.}
    \label{tab:register-umum}
    \small
    \begin{tabular}{@{}llllp{4.6cm}@{}}
        \toprule
        \textbf{Register} & \textbf{Byte tinggi} &
        \textbf{Byte rendah} & \textbf{Ukuran} & \textbf{Peran khas} \\
        \midrule
        \reg{AX} & \reg{AH} & \reg{AL} & 16 bit (2 $\times$ 8 bit)
        & Akumulator: hasil hitungan, operand wajib \asm{MUL}/\asm{DIV},
        layanan DOS \\
        \addlinespace
        \reg{BX} & \reg{BH} & \reg{BL} & 16 bit (2 $\times$ 8 bit)
        & Basis: penunjuk alamat memori pada mode terindeks \\
        \addlinespace
        \reg{CX} & \reg{CH} & \reg{CL} & 16 bit (2 $\times$ 8 bit)
        & Pencacah: pengulangan \asm{LOOP}, operasi string, posisi bit
        \asm{SHL}/\asm{SHR} \\
        \addlinespace
        \reg{DX} & \reg{DH} & \reg{DL} & 16 bit (2 $\times$ 8 bit)
        & Data: pasangan \reg{AX} pada \asm{MUL}/\asm{DIV}, alamat port
        I/O dan layanan DOS \\
        \bottomrule
    \end{tabular}
\end{table}

Akses byte ini sangat menghemat memori dan langkah kerja. Untuk
menyimpan nilai 5, cukup satu instruksi:

\begin{lstlisting}[language=MASM, numbers=none]
MOV AL, 5       ; hanya byte rendah AX yang berubah (AH tetap)
\end{lstlisting}

Sebaliknya, \lstasm|MOV AX, 5| menulis kedua bagian sekaligus:
\reg{AH} menjadi \hex{00} dan \reg{AL} menjadi \hex{05}. Perbedaan
ini penting karena \lstasm|MOV AL, 5| \emph{tidak} mengosongkan
\reg{AH}. Bila sebelumnya \reg{AH} bernilai \hex{FF}, maka setelah
\lstasm|MOV AL, 5| register \reg{AX} bernilai \hex{FF05}, bukan
\hex{0005}.

\subsection{Mengapa Peran Setiap Register Berbeda}

Keempat register umum sebenarnya sama-sama dapat menampung nilai apa
pun. Meskipun demikian, konvensi pemakaiannya berbeda-beda. Perbedaan
ini bukan sekadar kebiasaan, melainkan akibat cara prosesor 8086
menyandikan instruksinya \cite{intel1979}. Beberapa instruksi
memberikan perlakuan khusus kepada register tertentu.

\begin{description}
    \item[\reg{AX} sebagai akumulator] Operasi \asm{MUL}, \asm{DIV},
          \asm{IN}, \asm{OUT}, dan sebagian besar layanan DOS
          mengandaikan operand berada di \reg{AX} atau \reg{AL}.
          Instruksi yang melibatkan akumulator juga memiliki bentuk
          mesin yang lebih pendek sehingga lebih cepat.

    \item[\reg{BX} sebagai basis] Hanya \reg{BX} yang boleh menjadi
          register basis di dalam tanda kurung, misalnya
          \lstasm|MOV AL, [BX]|. Register umum lain tidak dapat
          dipakai dengan cara itu. Aturan mode pengalamatan ini
          dibahas pada Bab~\ref{chap:mode-pengalamatan}.

    \item[\reg{CX} sebagai pencacah] Instruksi \asm{LOOP}, \asm{REP},
          dan operasi string memakai \reg{CX} sebagai penghitung
          otomatis. Instruksi \asm{SHL} dan \asm{SHR} juga memakai
          \reg{CL} sebagai jumlah pergeseran.

    \item[\reg{DX} sebagai pasangan data] Pada perkalian dan
          pembagian, \reg{DX} bersama \reg{AX} membentuk bilangan 32
          bit: \reg{DX} menyimpan bagian tinggi. Pada layanan DOS,
          \reg{DX} menyimpan alamat string atau nomor port.
\end{description}

Pemahaman mengenai peran khas ini memudahkan membaca program orang
lain. Ketika melihat \lstasm|MOV CX, 10|, programmer langsung
menduga bahwa instruksi itu menyiapkan pengulangan sepuluh kali.

\begin{infobox}
Register umum boleh dipakai di luar peran khasnya. Yang tidak boleh
adalah melanggar aturan khusus: \lstasm|MUL BX| selalu memakai
\reg{AX} sebagai operand tersembunyi, sedangkan
\lstasm|MOV AX, [CX]| tidak sah karena \reg{CX} bukan register basis
yang diizinkan.
\end{infobox}

\subsection{Register Segmen}

Empat register berikutnya, yaitu \reg{CS}, \reg{DS}, \reg{SS}, dan
\reg{ES}, disebut \term{register segmen}. Register segmen tidak
menyimpan data yang dihitung, melainkan menyimpan \emph{nomor
segmen} tempat data atau instruksi berada. Pada arsitektur 8086,
alamat memori dibentuk dari pasangan nomor segmen dan \eng{offset},
seperti yang telah dijelaskan pada Bab~\ref{chap:segmentasi}.

\begin{table}[H]
    \centering
    \caption{Register segmen dan perannya pada pengalamatan memori.}
    \label{tab:register-segmen}
    \small
    \begin{tabular}{@{}lp{4.6cm}p{6.0cm}@{}}
        \toprule
        \textbf{Register} & \textbf{Nama lengkap} &
        \textbf{Peran} \\
        \midrule
        \reg{CS} & \eng{Code Segment} & Menunjuk segmen tempat
        instruksi program berada; berpasangan dengan \reg{IP} \\
        \addlinespace
        \reg{DS} & \eng{Data Segment} & Menunjuk segmen tempat
        variabel berada; dipakai pada hampir semua akses data \\
        \addlinespace
        \reg{SS} & \eng{Stack Segment} & Menunjuk segmen tempat
        \eng{stack} berada; berpasangan dengan \reg{SP} dan \reg{BP} \\
        \addlinespace
        \reg{ES} & \eng{Extra Segment} & Segmen tambahan; dipakai
        sebagai segmen tujuan pada operasi string bersama \reg{DI} \\
        \bottomrule
    \end{tabular}
\end{table}

Setiap akses memori pada dasarnya bermain dengan pasangan register
segmen--offset berikut.

\begin{itemize}
    \item \reg{CS}:\reg{IP} --- alamat instruksi berikutnya yang
          akan diambil.
    \item \reg{DS}:\reg{SI}, \reg{DS}:\reg{BX}, atau \reg{DS} plus
          offset --- alamat data umum.
    \item \reg{SS}:\reg{SP} atau \reg{SS}:\reg{BP} --- alamat di
          dalam \eng{stack}.
    \item \reg{ES}:\reg{DI} --- alamat tujuan operasi string.
\end{itemize}

Karena itu, program \file{.EXE} wajib menginisialisasi \reg{DS} di
awal program. Tanpa inisialisasi tersebut, \reg{DS} masih berisi
nilai peninggalan DOS dan instruksi seperti
\lstasm|MOV AX, nilai| membaca alamat yang salah.

\begin{peringatanbox}
Dua kesalahan yang sering terjadi pada register segmen.
\begin{enumerate}
    \item Menulis nilai segera langsung ke register segmen, misalnya
          \lstasm|MOV DS, 1000h|. Instruksi ini \emph{tidak sah}.
          Nilai hanya boleh dipindahkan dari register umum:
          \lstasm|MOV AX, 1000h| lalu \lstasm|MOV DS, AX|.
    \item Mengubah \reg{DS} di tengah program tanpa menyimpan nilai
          lamanya. Setelah itu, seluruh variabel pada segmen data
          lama tidak dapat dijangkau lagi sampai \reg{DS}
          dikembalikan.
\end{enumerate}
\end{peringatanbox}
